\scanpage{82}
\mspage{5}

\[
T_{E/S,0}^{\otimes2}\simeq T_{\Sigma_{\underline J},t}
\tag*{(30)}
\]

où $t$ est la section de $\Sigma_J^*=\Sigma_J\setminus D_J$ sur $S$, image de la section nulle de $E$. Ceci posé, on définit un \uline{épinglage de Legendre} de $E/S$ comme étant formé d'un iso

\[
\deleted{\unsure{T}}\underbrace{\Delta_{E/S}}_{T_{E/S,0}}^{\deleted{\unsure{\otimes2}}}
\xrightarrow{\overset{\unsure{\kappa|k}}{\sim}}
\deleted{T}\underline O_{\deleted{\unsure{\Sigma_{\underline J}}}}^{\deleted{\unsure{\natural|\otimes}}}(1)^{\natural}_t\simeq
\tag*{(31)}
\]

tel que \deleted{il existe un} l'isom \inserted{$\kappa^{\otimes2}$} déduit par passage au carré tensoriel s'identifie, via (30) et (27), à l'isom

\[
T_{\Sigma_J,t}\simeq T_{\Sigma_J,t}\wedge_{\{\pm1\}}(\mathord{\mathrm{I}\mkern-7mu\mu}_4^*\wedge\underline\omega_{\underline J})
\tag*{(32)}
\]

déduit d'\unsure{une|un} section

\[
\left\{\begin{aligned}
&\alpha\in\Gamma(S,\mathord{\mathrm{I}\mkern-7mu\mu}_4^*\wedge\underline\omega_{\underline J})\qquad\text{i.e d'un \uline{iso}}\\
&\alpha:\mathord{\mathrm{I}\mkern-7mu\mu}_4^*\simeq\underline\omega_{\underline J}
\end{aligned}\right.
\tag*{(\unsure{3|8}3)}
\]

(qui est déterminé de façon unique en termes
% archive sheet number: 81
